\section{Few negative directions: quadratic and mixed-integer quadratic problems}\label{sec:qp}

This section treats the first structural route, in which the nonconvexity of
the objective is confined to a small number $k$ of directions. For a
quadratic objective on a bounded polytope or mixed polytope, $k$ is the
negative inertia of the Hessian. A rational factorization
$A=P-\alpha T^{\mathsf T} T$ with $P\succeq0$ moves the nonconvex part to an
auxiliary space $\R^k$, where the search of Section~\ref{sec:count} runs. The
remaining recourse is a convex quadratic program, a convex mixed-integer
quadratic program, or a separable scalar problem, and it is solved exactly at
every query.

The principal result is Theorem~\ref{thm:qp:gauss}. Under a finite
Gaussian-like law on all original linear coefficients, selected from the
base data, it computes an exact optimizer of the sampled quadratic or
mixed-integer quadratic program on every draw. Its expected bit work is
$f(n_z,k,1+\nu\diam(\mathcal P)/\sigma)\,(I+1)^{C}$ with an absolute
exponent $C$, where $n_z$ is the number of integer variables, $\nu$ the
magnitude of the most negative eigenvalue of the Hessian, and $\mathcal P$
the polyhedral relaxation. The
other results are variants that differ in the perturbation law or the
recourse structure: noise aligned with the negative directions
(Theorem~\ref{thm:qp:aligned}), uniform noise on all coefficients, whose
bound contains a power of the dimension that grows with $k$
(Theorem~\ref{thm:qp:uniform}), at most two negative directions with a bound
linear in a numerical ratio (Theorem~\ref{thm:qp:two}), and separable mixed
recourse with arbitrarily many integer variables
(Theorems~\ref{thm:qp:sep} and~\ref{thm:qp:sep-gauss}).
Table~\ref{tab:qp:regimes} summarizes them. Every result concerns the exact
optimum of the sampled objective. Proofs are in Appendix~\ref{app:qp}.

\subsection{Instances, exact primitives, and factorizations}\label{sec:qp:model}

\begin{definition}[Quadratic instances]\label{def:qp:instance}
A \emph{base instance} consists of a symmetric matrix $A\in\Q^{n\times n}$,
a vector $b\in\Q^n$, a matrix $D\in\Q^{r\times n}$ and a vector $e\in\Q^r$
such that $\mathcal P=\{x:Dx\le e\}$ is nonempty and bounded, a number
$n_z\in\{0,\dots,n\}$ of integer variables, and a rational noise width
$\sigma>0$. The feasible set is
$X=\mathcal P\cap(\Z^{n_z}\times\R^{n-n_z})$, and the objective is
$F(x)=\frac12x^{\mathsf T} Ax+b^{\mathsf T} x$. We assume $X\ne\emptyset$.
If this is not promised, one convex-MIQP feasibility solve detects
infeasibility before any optimum, growth threshold, or search is defined.
Lower-dimensional polytopes and redundant
rows are allowed. Let $n_-(A)$ be the number of negative eigenvalues of $A$,
counted with multiplicity, and put $\nu=\max\{0,-\lambda_{\min}(A)\}$. A
\emph{perturbation} adds a linear term: \emph{ambient} noise gives
$F_\gamma(x)=F(x)+\gamma^{\mathsf T} x$ with $\gamma\in\R^n$, and noise
\emph{aligned} with a factor $T\in\Q^{k\times n}$ gives
$F(x)+\xi^{\mathsf T} Tx$ with $\xi\in\R^k$; these are the ambient and aligned
models of Definition~\ref{def:model:perturbation}. We write $f^*$ for the
minimum of the sampled objective over $X$.
\end{definition}

Aligned noise perturbs the original coefficients by $T^{\mathsf T}\xi$. These
coefficients are correlated and lie in a $k$-dimensional subspace, so aligned
noise is not independent noise on the original coefficients. Both models are
treated because they lead to different bounds.

\begin{proposition}[Exact convex primitives]\label{prop:qp:primitives}
\begin{enumerate}[label=(\alph*)]
\item A convex quadratic program $\min\{\frac12x^{\mathsf T} Px+p^{\mathsf T} x:Dx\le e\}$
 with $P\succeq0$ and rational data of bit length $L_{\rm in}$ is solved
 exactly in $\poly(L_{\rm in})$ bit operations: infeasibility and
 unboundedness are detected, and otherwise the optimal value and a rational
 optimal solution are returned
 \cite{kozlov1980-the-polynomial-solvability-of-convex}.
\item A convex quadratic program with $n_z$ integer variables and the other
 data as in (a) is solved exactly, with the same outputs, in
 $f(n_z)L_{\rm in}^{c_0}$ bit operations for a function $f$ and an absolute
 constant $c_0$ \cite[arXiv v2, Theorem~3]{pia2025-convex-quadratic-sets-and-the}.
\item Rational linear programs are solved exactly in polynomial time; in
 particular a vertex of a nonempty rational polytope can be found
 \cite{grotschel1988-geometric-algorithms-and-combinatorial-optimization}.
\end{enumerate}
\end{proposition}

When the algorithm of (b) returns an optimal solution, we fix its integer
part, which lies in the box of linear-programming bounds and therefore has
polynomial bit length, and re-solve the continuous slice by (a). This
\emph{polishing} returns an optimal solution of polynomial bit length, so
that the possibly large factor $f(n_z)$ never enters the size of later
instances. The factor $f(n_z)$ is absent when $n_z=0$.

\begin{definition}[Convexifying factorization]\label{def:qp:factor}
A \emph{convexifying factorization} of $A$ is a rational number $\alpha>0$
and a rational matrix $T\in\Q^{k\times n}$ with $P=A+\alpha T^{\mathsf T} T\succeq0$.
It has a supplied rational \emph{frame constant} $c_{\rm fr}\in(0,1]$, counted
in $I$, if
$c_{\rm fr}I_k\preceq TT^{\mathsf T}\preceq I_k$. The number $k$ of rows satisfies
$k\ge n_-(A)$, with equality for the factorization of
Lemma~\ref{lem:qp:normalize}; in theorems with a supplied factorization, $k$
denotes its number of rows.
\end{definition}

Positive semidefiniteness of $P$ and a supplied frame constant are verified
exactly by symmetric elimination, so a supplied factorization is a checked
certificate rather than a promise. No relation between $\ker P$ and $\ker T$
is required.

\begin{lemma}[Intrinsic rational normalization]\label{lem:qp:normalize}
There is a deterministic algorithm, polynomial in the bit length of $A$, that
either certifies $A\succeq0$, or computes rational $\beta$ with
$\nu\le\beta<2\nu$ and a rational matrix $T\in\Q^{k\times n}$, with
$k=n_-(A)$, such that
\[
 A+2\beta\,T^{\mathsf T} T\succeq0,
 \qquad
 \tfrac{63}{64}I_k\preceq TT^{\mathsf T}\preceq I_k,
\]
and the rows of $T$ lie in the range of $A$. Thus $\alpha=2\beta$ gives a
convexifying factorization with $2\nu\le\alpha<4\nu$ and frame constant
$63/64$. It uses no real eigenvectors, no bound on $\nu$, and no lower bound
on nonzero eigenvalues as input.
\end{lemma}

The construction combines exact rational Jacobi rotations, which give an
orthogonal near-diagonalization \cite{pia2026-rational-jacobi-rotations-and-the},
with an exact projection onto the range of $A$. The projection is essential.
For $A=\diag(-1,0)$ and a unit vector $t=(\cos\theta,\sin\theta)$ with
$\sin\theta\ne0$, the determinant of $A+2tt^{\mathsf T}$ is $-2\sin^2\theta<0$, so an
approximate negative eigenvector with a component in $\ker A$ does not give a
convexifier, however small that component is. For the factorization of the lemma,
$P=A+2\beta T^{\mathsf T}T$ is positive definite on the range of $A$ and
vanishes on $\ker A$, so $\ker P=\ker A\subseteq\ker T$. This inclusion
holds automatically, but none of the proofs below uses it.

\subsection{The auxiliary value function}\label{sec:qp:aux}

Let $X\subset\R^n$ be nonempty and compact, let $F:X\to\R$ be continuous,
let $T\in\R^{k\times n}$, and let $\Lambda=\diag(\lambda_1,\dots,\lambda_k)$
be positive definite; in all results except Theorem~\ref{thm:qp:sep-gauss},
$\Lambda=\alpha I$. For a \emph{residual} $\zeta\in\R^n$ and a \emph{factor
coefficient} $\xi\in\R^k$ define
\[
 W_\zeta(a)=\min_{x\in X}\Bigl[F(x)+\zeta^{\mathsf T} x+\tfrac12(a-Tx)^{\mathsf T}\Lambda(a-Tx)\Bigr],
 \qquad
 V(a)=W_\zeta(a)+\xi^{\mathsf T} a,
\]
and put $\gamma=T^{\mathsf T}\xi+\zeta$. A point attaining the minimum defining
$W_\zeta(a)$ is an \emph{attaining witness} at $a$. For a quadratic $F$ with a
convexifying factorization and $\Lambda=\alpha I$, the inner problem has the
convex objective $\frac12x^{\mathsf T} Px+(b+\zeta-\alpha T^{\mathsf T} a)^{\mathsf T} x$ up to a
constant, so $W_\zeta(a)$ is one convex QP or convex MIQP.

\begin{lemma}[Auxiliary value function]\label{lem:qp:aux}
In this setting:
\begin{enumerate}[label=(\alph*)]
\item \emph{(Square completion.)} For every $a$,
 \[
  V(a)=\min_{x\in X}\Bigl[F_\gamma(x)+\tfrac12(a-Tx+\Lambda^{-1}\xi)^{\mathsf T}\Lambda(a-Tx+\Lambda^{-1}\xi)\Bigr]
  -\tfrac12\xi^{\mathsf T}\Lambda^{-1}\xi .
 \]
 Consequently $V^*=\inf_{\R^k}V=f^*-\frac12\xi^{\mathsf T}\Lambda^{-1}\xi$, attained at
 $Tx^*-\Lambda^{-1}\xi$ for every minimizer $x^*$ of $F_\gamma$, and every
 attaining witness $x_a$ at any $a$ satisfies $F_\gamma(x_a)-f^*\le V(a)-V^*$.
\item \emph{(Upper curvature.)} $W_\zeta(a)-\frac12a^{\mathsf T}\Lambda a$ is concave,
 so $W_\zeta$ and $V$ have upper coordinate curvature at most
 $(\lambda_1,\dots,\lambda_k)$.
\item \emph{(Every attaining witness gives a small active vector.)} For
 every $a$ and every attaining witness $x_a$, the vector
 $g=\Lambda(a-Tx_a)+\xi$ satisfies
 \[
  V(z)\le V(a)+g^{\mathsf T}(z-a)+\tfrac12(z-a)^{\mathsf T}\Lambda(z-a)\quad(z\in\R^k),
  \qquad
  g^{\mathsf T}\Lambda^{-1}g\le2\bigl(V(a)-V^*\bigr).
 \]
\item \emph{(A fixed box.)} If $\ell_i\le(Tx)_i\le u_i$ for all $x\in X$ and
 $|\xi_i|\le s_i$, then
 $\mathcal A=\prod_i[\ell_i-s_i/\lambda_i,\,u_i+s_i/\lambda_i]$ contains a
 global minimizer of $V$, and $\min_{\mathcal A}V=V^*$.
\end{enumerate}
\end{lemma}

Part (c) is the inequality that replaces differentiability of $W_\zeta$
throughout this section. It holds for \emph{every} attaining witness,
including witnesses at integer ties, on flat optimal faces, and at points
where $W_\zeta$ is not differentiable. No derivative of $W_\zeta$ and no
selection of a favorable subgradient is needed. For ambient noise and a
factor of full row rank, the decomposition is
\[
 \Psi=(TT^{\mathsf T})^{-1}T,\qquad \xi=\Psi\gamma,\qquad
 \zeta=\Pi_T\gamma,\qquad \Pi_T=I-T^{\mathsf T}\Psi,
\]
so that $\gamma=T^{\mathsf T}\xi+\zeta$ and $T\zeta=0$. The coordinates of $\xi$ are
in general dependent, and $\xi$ and $\zeta$ are dependent under uniform
noise; Section~\ref{sec:qp:ambient} shows why this matters.

\subsection{The inverse-growth route for at most two negative directions}\label{sec:qp:growth}

When the sampled problem has positive point growth, the auxiliary problem
inherits growth, and a search without closure terminates exactly by rational
reconstruction. This deterministic result is combined with the growth tail.

\begin{theorem}[Growth-conditioned exact algorithm]\label{thm:qp:conditioned}
There is a deterministic algorithm for continuous instances ($n_z=0$) with
$k=n_-(A)\ge1$ and rational data of bit length $L_{\rm in}$ with the
following properties.
\begin{enumerate}[label=(\alph*)]
\item Every output is an exact rational global minimizer together with the
 exact optimal value; this holds on every instance.
\item If the minimizer $x^*$ is unique and $F(x)-F(x^*)\ge g\norm{x-x^*}^2$ on
 $X$ for some $g>0$, the algorithm halts. With $\kappa=2+4\nu/g$, it
 processes at most $4^k(2\sqrt{k\kappa}+2)^k$ cells per refinement level,
 stops after at most $J_g\le\poly(L_{\rm in})+\log_2(1+\kappa)$ levels, and
 uses at most $c_k(1+\sqrt\kappa)^k(J_g+1)\poly(L_{\rm in}+J_g)$ bit
 operations, where $c_k$ depends only on $k$.
\item For $k\le2$ and $Z=\max\{1,\nu/g\}$ the bound in (b) is at most
 $\poly(L_{\rm in})\,Z^{k/2}(1+\log Z)^{d_0}$, where $d_0$ and the polynomial
 are absolute.
\end{enumerate}
The algorithm needs neither $g$ nor $\nu$.
\end{theorem}

The algorithm normalizes $A$ by Lemma~\ref{lem:qp:normalize}, runs the
corrected-corner search of Definition~\ref{def:count:search} on $W_0$ over the
box of ranges of $Tx$, and attempts two rational reconstructions at every
level: of the optimal value from the certified interval of
Theorem~\ref{thm:count:cells}(v), and of the common auxiliary minimizer $Tx^*$
from the best corner. It accepts only a feasible witness whose value equals
the reconstructed optimal value, so a premature reconstruction is harmless.
Growth enters only the termination bound: by Lemma~\ref{lem:qp:aux} and the
triangle inequality, $W_0(a)-f^*\ge g_W\norm{a-Tx^*}^2$ with
$\alpha/g_W\le\kappa$. Every retained cell has a near-optimal corner in the
ball of radius $h_j\sqrt{k\kappa}/2$ around $Tx^*$, and the whole cell lies
in that ball enlarged by its diameter.

\begin{corollary}[Growth-conditioned mixed-integer algorithm]\label{cor:qp:conditioned-mixed}
Theorem~\ref{thm:qp:conditioned} also holds for nonempty mixed feasible sets
of Definition~\ref{def:qp:instance}, with an additional factor $f(n_z)$
in its bit-work bounds, where $f$ is the function of
Proposition~\ref{prop:qp:primitives}(b). A feasible rational optimizer and
its exact value are returned whenever the algorithm halts; positive point
growth guarantees termination. The polynomial exponents are independent
of $n_z$ and $k$, and the factor $c_k$ remains dependent only on $k$.
No growth modulus is needed as input.
\end{corollary}

To use Theorem~\ref{thm:qp:conditioned} under a finite law, the growth tail
must be transferred to that law with an atom term that does not depend on the
threshold.

\begin{lemma}[Explicit growth sections for quadratics]\label{lem:qp:growth-sections}
Let $F$ be quadratic and let $X=\mathcal P\cap(\Z^{n_z}\times\R^{n-n_z})$ as in
Definition~\ref{def:qp:instance}, and let $\mathsf F=Z\cdot2^r$, where $Z$ is
the number of integer points in the box of linear-programming bounds of the
integer variables ($Z=1$ if $n_z=0$). For every $\varepsilon>0$, every
coordinate $i$, and all fixed values of the other coefficients, the set of
values of $\gamma_i$ for which $g_*(\gamma)<\varepsilon$ is a union of at most
$8(\mathsf F+1)^2$ intervals. Consequently, with $S=\sum_iw_i$ the sum of the
coordinate widths of $X$,
\[
 \Prob\{g_*<\varepsilon\}\le\frac{S\varepsilon}\sigma+\frac{16n(\mathsf F+1)^2}M
 \quad\text{and}\quad
 \Prob\{g_*<\varepsilon\}\le\sqrt{\tfrac2\pi}\,\frac{S\varepsilon}\sigma+16n(\mathsf F+1)^22^{-b}
\]
for independent coordinates with law $U_{\sigma,M}$ and
$\mathcal G_{b,\sigma}$, respectively, for every $\varepsilon>0$.
\end{lemma}

\begin{theorem}[Expected exact work for at most two negative directions]\label{thm:qp:two}
Let the base instance be continuous ($n_z=0$) with $k=n_-(A)\le2$, let
$S$ be the sum of the coordinate widths of $\mathcal P$, and put
$B=\max\{2,2^r\}$ and $\mathsf D=8(2^r+1)^2$.
\begin{enumerate}[label=(\alph*)]
\item Let $M$ be the least power of two with $M\ge2n\mathsf DB$, so
 $\log_2M=O(r+\log(n+1))$, and let $\gamma_1,\dots,\gamma_n$ be independent
 with law $U_{\sigma,M}$. There is an algorithm that returns an exact
 rational minimizer of $F_\gamma$ and the exact optimal value on every draw,
 with expected bit work at most $(I+1)^{C}(1+\nu S/\sigma)$ for an absolute
 constant $C$.
\item The same holds for independent coordinates with law
 $\mathcal G_{b,\sigma}$, where $b$ is the least positive integer satisfying
 $2^b\ge2n\mathsf DB$, hence $b=O(r+\log(n+1))$.
\end{enumerate}
Either law may use a larger resolution with $\log M\le\poly(I)$ or
$b\le\poly(I)$, respectively; the bit polynomial then includes that
resolution. This permits a common law with a closure theorem while
charging its sampling and coefficient lengths.
The algorithm needs neither $\nu$, $S$, nor the growth of the sampled
instance. For $k=0$ one exact convex QP solve suffices.
\end{theorem}

The algorithm interleaves, bit operation by bit operation, the algorithm of
Theorem~\ref{thm:qp:conditioned} and the exact fallback of
Lemma~\ref{lem:qp:face} on the same draw, and stops when either finishes. The
fallback costs at most $B\poly(L_{\rm in})$ on every draw. With
$p=k/2$ and $Z=\max\{1,\nu/g_*\}$ ($Z=\infty$ when $g_*=0$), the work is at
most $\poly(L_{\rm in})\min\{B,Z^p(1+\log Z)^{d_0}\}\le
\poly(L_{\rm in})(1+p^{-1}\log B)^{d_0}\min\{B,Z^p\}$. The growth tail with a
residual $\beta\le1/B$ gives $\Prob\{Z>t\}\le r_0/t+\beta$ for $t\ge1$, where
$r_0=\nu S/\sigma$, and therefore
\[
 \E\min\{B,Z^p\}\le1+r_0\int_1^Bt^{-1/p}\,dt+\beta B
 \le\begin{cases}2+r_0,&k=1,\\ 2+r_0\log B,&k=2.\end{cases}
\]
All atoms and all draws with $g_*=0$ are included in the term $\beta B$.

\begin{remark}[Why the moment argument stops at two directions]\label{rem:qp:moments}
For $k\ge3$ the same integral is
$r_0(B^{1-2/k}-1)/(1-2/k)$, which is exponential in $r$. This is a limitation
of the argument, not a lower bound for any algorithm. It cannot be repaired by
a sharper tail. Let $X=[-w_0/2,w_0/2]\times Y$ and $F(x_0,y)=G(y)$, where $G$
is a quadratic with $k$ negative eigenvalues of magnitude at most $\nu$ on a
polytope $Y$. For a continuous uniform perturbation, moving $x_0$ to the other
endpoint shows $g_*\le|\gamma_0|/w_0$. Hence
$\Prob\{Z>t\}\ge\nu w_0/(\sigma t)$ for $t\ge\max\{1,\nu w_0/\sigma\}$, and
$\E\min\{B,Z^{k/2}\}\ge a_0\int_{t_0}^Bs^{-2/k}\,ds$ when $B>t_0$, with
$a_0=\nu w_0/\sigma$ and $t_0=\max\{1,a_0^{k/2}\}$. For fixed $a_0>0$ and
sufficiently large $B$, this is of order
$(\nu w_0/\sigma)B^{1-2/k}$ for $k\ge3$ and large $B$. The inverse-growth
power $Z^{k/2}$ is therefore not integrable to a polynomial bound beyond two
negative directions, even though the tail of
Theorem~\ref{thm:count:growth-tail} is sharp. The closure method below does
not integrate an inverse growth power and has no such restriction. The two
results are not comparable: Theorem~\ref{thm:qp:two} is linear in
$\nu S/\sigma$, while the closure bounds are of degree $k$ in
$\nu\diam(\mathcal P)/\sigma$. They can be applied to the same sampled
objective: part (b) of Theorem~\ref{thm:qp:two} holds for the law of
Theorem~\ref{thm:qp:gauss} once its accuracy also satisfies
$2^b\ge2n\mathsf DB$, a requirement of polynomial bit length that does not
depend on the support.
\end{remark}

\subsection{Exact closure of cells}\label{sec:qp:closure}

\begin{lemma}[Smallest-face candidates and the quadratic fallback]\label{lem:qp:face}
Let $G(x)=\frac12x^{\mathsf T} Hx+p^{\mathsf T} x$ be a rational quadratic, with $H$
symmetric of any inertia, and let $\mathcal Q=\{x:Dx\le e\}$ be a nonempty
bounded rational polytope.
\begin{enumerate}[label=(\alph*)]
\item There are a global minimizer $\hat x$ of $G$ on $\mathcal Q$ and a set
 $S$ of linearly independent rows active at $\hat x$ such that
 $K_S=\bigl(\begin{smallmatrix}H&D_S^{\mathsf T}\\D_S&0\end{smallmatrix}\bigr)$ is
 nonsingular and $\hat x$ is the $x$-part of $K_S^{-1}(-p;e_S)$.
\item Enumerating all row subsets, solving the nonsingular systems, keeping
 the feasible solutions, and returning one of least value finds an exact
 rational global minimizer in $2^r\poly(L_{\rm in})$ bit operations, on every
 instance. Multiplier signs and definiteness tests are unnecessary.
\item For a mixed polytope, enumerating the $Z$ integer points of the box of
 linear-programming bounds and applying (b) to every nonempty slice is exact
 in $B\poly(L_{\rm in})$ bit operations with $B=\max\{2,Z2^r\}$.
\item For a box in $\R^k$, the $3^k$ faces (each coordinate at its lower
 bound, at its upper bound, or free) suffice in (b).
\end{enumerate}
In particular, every rational QP and MIQP on a nonempty bounded mixed polytope has a
rational global minimizer, and an optimal value, of polynomial bit length.
\end{lemma}

The proof chooses a global minimizer on a face of smallest dimension. Its
tangent Hessian is positive semidefinite; a null tangent direction would keep
the objective constant up to the boundary of the face and produce a minimizer
on a smaller face. This argument handles singular Hessians, continua of
minimizers, and lower-dimensional polytopes, but it is specific to quadratics.

Let a base instance and a convexifying factorization be fixed, with
$\Lambda=\alpha I$. A \emph{label} is the integer part $z$ of a point of $X$,
and the \emph{slice} of $z$ is the set of continuous parts $x_c$ with
$(z,x_c)\in X$. Write $W_\zeta^z$ for the function $W_\zeta$ with $X$
replaced by the points of label $z$, so that $W_\zeta=\min_zW_\zeta^z$.

\begin{lemma}[Critical pieces of the convex recourse]\label{lem:qp:pieces}
For every feasible label $z$ and every rational query $(a,\zeta)$, polynomial-time
rational computations after one exact solve of the slice problem return a
set $S$ of linearly independent rows of the slice constraints, active at an
optimal slice point, such that the following hold.
\begin{enumerate}[label=(\alph*)]
\item The stationarity system of the slice with active rows $S$ is
 nonsingular. Its solution is affine in $(a',\zeta')$:
 $x_{z,S}(a',\zeta')=\Xi_Sa'+Y_S\zeta'+x^0_{z,S}$, with multipliers
 $\lambda_{z,S}(a',\zeta')$ affine as well.
\item The \emph{region}
 $\mathcal R_{z,S}(\zeta')=\{a':x_{z,S}(a',\zeta')\text{ feasible},\
 \lambda_{z,S}(a',\zeta')\ge0\}$ contains $a$ when $\zeta'=\zeta$.
 It is defined by at most $r+n$
 affine inequalities whose normals do not depend on $\zeta'$ and whose
 offsets are affine in $\zeta'$. For every $a'\in\mathcal R_{z,S}(\zeta')$,
 the point $x_{z,S}(a',\zeta')$ attains $W^z_{\zeta'}(a')$.
\item The quadratic
 $q_{z,S}(a',\zeta')=F(x_{z,S})+\zeta'^{\mathsf T} x_{z,S}+\frac\alpha2\norm{a'-Tx_{z,S}}^2$
 equals $W^z_{\zeta'}$ on the region and satisfies the identity
 $\nabla_{a'}q_{z,S}(a',\zeta')=\alpha\bigl(a'-Tx_{z,S}(a',\zeta')\bigr)$ for
 all $(a',\zeta')$. Its Hessian $H_S=\alpha(I-T\Xi_S)$ does not depend on
 $\zeta'$, $z$, $b$, or $e$, and $\norm{H_S}\le\Theta$, where
 $\Theta=\alpha(1+kn\tau_TU_0)$, $\tau_T=\max_{il}|T_{il}|$, and
 $U_0=(2n)!\,C_0^{2n}$ with $C_0$ the largest absolute value of the integers
 obtained by clearing a common denominator of the entries of $P$, $D$, and
 $\alpha T^{\mathsf T}$.
\end{enumerate}
There are at most $R=Z\cdot2^r$ possible pairs $(z,S)$.
For every feasible fixed label these regions cover all real parameter
pairs $(a,\zeta)$; only extraction at rational queries is computational.
The region description does not require full dimension, an irredundant
description, or a relation between $\ker P$ and $\ker T$.
\end{lemma}

The gradient identity in (c) is algebraic. It holds on all of $\R^k$, even
when the region has empty interior; it is not obtained by differentiating
$W_\zeta$ along the region. At the query, $x_{z,S}(a,\zeta)$ is an attaining
witness of $W^z_\zeta(a)$, so Lemma~\ref{lem:qp:aux}(c) applies to the
vector $\nabla q_{z,S}(a,\zeta)+\xi$ whenever $z$ is a winning label. The bound
$\Theta$ uses only the matrices of the response to $a'$. The sampled
coefficients enter only offsets, so $\Theta$ is fixed before sampling. Its
numerical size is large, but it enters the algorithm only through its
logarithm.

Under the usual uniqueness and regularity assumptions, affine optimizer
formulas on polyhedral critical regions, with quadratic value formulas, are
standard in multiparametric quadratic programming
\cite{bemporad2002-the-explicit-linear-quadratic-regulator,tndel2003-an-algorithm-for-multi-parametric,patrinos2011-convex-parametric-piecewise-quadratic-optimization},
and the reduction of structured indefinite quadratic programs to the global
minimization of a parametric convex value function, with recovery of
optimizers from the minimizing parameter, appears in Ding's thesis
\cite{ding1996-a-parametric-solution-for-local}. Lemma~\ref{lem:qp:pieces}
is the form of this machinery used here: one region per query, extracted in
polynomial time without enumerating regions, with a Hessian bound that is
independent of the sampled coefficients. What the expected-work theorems add
is the analysis of the cells that a region does not close.

For a continuous instance a region certifies the value of $W_\zeta$ on a
whole cell. For integer variables it certifies only the value of one label,
and a second certificate is needed.

\begin{lemma}[Best-other-label gaps]\label{lem:qp:gap}
Let $[\ell_i,u_i]$ be the range of $(Tx)_i$ over $\mathcal P$, and put
$\bar w=\max\{1,\sum_i(u_i-\ell_i)\}$. At a query $v$, let $z_*$ be a winning
label and put $\Delta(v)=\min_{z\ne z_*}W^z_\zeta(v)-W_\zeta(v)$, with
$\Delta(v)=+\infty$ if no other label is feasible.
\begin{enumerate}[label=(\alph*)]
\item $\Delta(v)$ and a witness of the best other label are computed by $2n_z$
 additional exact convex MIQP solves, each with one extra inequality
 $x_i\le z_{*,i}-1$ or $x_i\ge z_{*,i}+1$.
\item Let $C$ be a cell with corner $v$ and side lengths at most $h$. If
 $C\subseteq\mathcal R_{z_*,S}(\zeta)$ and $\Delta(v)\ge\theta_{\rm gap}h$
 with $\theta_{\rm gap}=\alpha k\bar w$, then $W_\zeta=q_{z_*,S}(\cdot,\zeta)$
 on $C$, and $x_{z_*,S}(a,\zeta)$ attains $W_\zeta(a)$ for every $a\in C$.
\item If $V(v)-V^*\le2E$ with $E\le k\alpha h^2/8$ and $h\le1$, and the gap
 condition of (b) fails, then two different labels have feasible points with
 $F_\gamma\le f^*+C_{\rm gap}h$, where $C_{\rm gap}=k\alpha/4+\theta_{\rm gap}$.
\end{enumerate}
\end{lemma}

Equality in the gap test is allowed: another label may tie inside the cell
without changing the certified value or the feasible witness. Equal winning
labels at all corners do not suffice (Example~\ref{ex:qp:corner-labels}).

\begin{lemma}[Integer-label isolation]\label{lem:qp:isolation}
Let $\mathcal Z$ be a nonempty finite set of integer vectors in
$\prod_{i=1}^{n_z}[L_i,U_i]$, let $h:\mathcal Z\to\R$ be arbitrary, and let
$\gamma\in\R^{n_z}$ have independent coordinates with interval concentrations
$\pi_i$. Let $\Delta$ be the difference between the two smallest values of
$h(z)+\gamma^{\mathsf T} z$ at distinct labels ($\Delta=+\infty$ if
$|\mathcal Z|=1$). Then for every $\varepsilon\ge0$
\[
 \Prob\{\Delta\le\varepsilon\}\le\sum_{i=1}^{n_z}(U_i-L_i)\,\pi_i(2\varepsilon).
\]
\end{lemma}

In the application, $h(z)$ is the optimal value of the slice of $z$ with its
continuous noise, conditioned on that noise, so the continuous recourse may
be arbitrary. The widths $U_i-L_i$ can be exponentially large, but they enter
the algorithms only through logarithms. Isolation of the best discrete
solution under random objective coefficients is classical
\cite{beier2006-typical-properties-of-winners-and,roglin2007-smoothed-analysis-of-integer-programming};
the lemma is the version used here, with finite laws and arbitrary label
costs.

\begin{example}[Equal corner winners do not certify a label]\label{ex:qp:corner-labels}
Let $z\in[0,1]$ be continuous and $y\in\{0,1\}$ integer, and let
$F(z,y)=\frac12zy+\frac78y^2-\frac{15}{16}y$, $\alpha=1$, $T=(1,1/2)$. Then
$P=\bigl(\begin{smallmatrix}1&1\\1&2\end{smallmatrix}\bigr)\succ0$. On
$a\in[0,1]$ the label $y=0$ has slice value $0$ and the label $y=1$ has slice
value $a^2/2-a/2+1/16$. The label $y=0$ wins strictly at both endpoints, and
its slice formula is valid on the whole cell, but $y=1$ wins at $a=1/2$. The
gap test of Lemma~\ref{lem:qp:gap} refuses this cell; a test based on equal
corner labels would accept it and return a wrong value. Rescaling to
$\alpha=4$, $T=(1/2,1/4)$, and the cell $[0,1/2]$ gives the same example with
$\norm T<1$.
\end{example}

\begin{definition}[Closure search]\label{def:qp:closure-search}
Let a base instance, a convexifying factorization, a residual $\zeta$, a
factor coefficient $\xi$, a box $\mathcal A$ as in Lemma~\ref{lem:qp:aux}(d),
and a level cap $J$ be given. Run the corrected-corner search of
Definition~\ref{def:count:search} on $V$ over $\mathcal A$ with isotropic
meshes and curvature $\alpha$. At every corner $v$ of a processed cell,
compute $W_\zeta(v)$ exactly with a polished attaining witness, a winning
label $z_v$, a piece $(z_v,S_v)$ of Lemma~\ref{lem:qp:pieces}, and, if
$n_z\ge1$, the gap $\Delta(v)$ of Lemma~\ref{lem:qp:gap}. A processed cell
$C$ of level $j$ is \emph{closed} if for some corner $v$ all corners of $C$
lie in $\mathcal R_{z_v,S_v}(\zeta)$ and, when $n_z\ge1$,
$\Delta(v)\ge\theta_{\rm gap}h_j$. A closed cell is minimized exactly as the
quadratic $q_{z_v,S_v}(\cdot,\zeta)+\xi^{\mathsf T} a$ on $C$ by
Lemma~\ref{lem:qp:face}(d), and the minimizer's affine witness is stored. If
no unclosed cell is retained at some level $j\le J$, return the witness of the
incumbent. Otherwise, after level $J$, run the exact fallback of
Lemma~\ref{lem:qp:face}(c) on the same sampled objective.
\end{definition}

By Lemmas~\ref{lem:qp:pieces} and~\ref{lem:qp:gap}(b), every closure step is
exact, so Theorem~\ref{thm:count:cells} applies to $V$ on $\mathcal A$. When
the search stops before the fallback, its incumbent attains $V^*$, and by
Lemma~\ref{lem:qp:aux}(a) the stored witness is a global minimizer of the
sampled objective. The output is therefore exact on every draw, with a
rational minimizer and value of polynomial bit length. The search trace,
consisting of closed cells with their pieces and discarded cells with their
corner bounds, is a verifiable global certificate; for $n_z\ge1$ its
verification includes rerunning the exact mixed-integer solves.

The probabilistic analysis rests on one deterministic fact: a retained
unclosed cell forces the noise close to one of finitely many hyperplanes that
are fixed before sampling.

\begin{lemma}[Unresolved cells force fixed tubes]\label{lem:qp:tube}
There is a family of at most $K=R(r+n+1)$ affine functionals
$\ell_\kappa(\xi,\zeta)=u_\kappa^{\mathsf T}\xi+w_\kappa^{\mathsf T}\zeta+c_\kappa$ with
$\norm{u_\kappa}=1$, determined by the base data and the factorization, with
the following property. Suppose that a cell of level $j$ is retained and not
closed, that its corner $v$ satisfies $V(v)-V^*\le2E_j$, that $z_v$ wins at
$v$, and that some corner of the cell lies outside
$\mathcal R_{z_v,S_v}(\zeta)$. Then
$|\ell_\kappa(\xi,\zeta)|\le\sqrt k(\alpha+\Theta)h_j$ for some $\kappa$. If
$T$ has full row rank and $\norm T\le1$, then
$\ell_\kappa(\Psi\gamma,\Pi_T\gamma)=L_\kappa^{\mathsf T}\gamma+c_\kappa$ with a
fixed vector $L_\kappa$ of norm at least $1$.
\end{lemma}

If the piece matrix $H_S$ is invertible, the region boundary crossed inside
the cell is mapped by the affine map $a\mapsto-\nabla_aq_{z,S}(a,\zeta)$ to a
hyperplane in the space of $\xi$; by Lemma~\ref{lem:qp:aux}(c), $\xi$ lies
close to the image of $v$. If $H_S$ is singular, the whole image of that
map lies in a hyperplane. The normals are fixed, and only the offsets depend on
$\zeta$, affinely. A region row with zero normal cannot be violated in a cell
that contains the query, so no exceptional case arises. The family is a proof
device; the algorithm never constructs it.

\begin{example}[An atom that requires the fallback]\label{ex:qp:atom}
Let $X=[0,1]^2$, $F(x,y)=xy-(x-y)/2$, $T=(1/2,-1/2)$, $\alpha=4$, and
$\sigma=1$, so $P=I$. Under aligned noise the endpoint draw $\xi=1$, which
belongs to every grid $U_{1,M}$, gives the sampled objective $xy$.
A direct computation gives $V(a)=-1/8$ for all $a\in[-3/4,1/4]$, with two
critical pieces, $\{x=1/2+2a,\ y=0\}$ for $a\ge-1/4$ and
$\{x=0,\ y=-1/2-2a\}$ for $a\le-1/4$. The fixed box is $[-3/4,3/4]$, and the
switch $a=-1/4$ lies at relative position $1/3$, which is never a dyadic grid
node. The cell containing $-1/4$ is therefore never closed, and its corners are
globally optimal, so it is retained at every level. The draw has probability
$1/M$ and is solved by the fallback. Closure alone does not make the
algorithm terminate on this draw; the fallback is necessary.
\end{example}

\subsection{Main theorems}\label{sec:qp:main}

In the next three theorems, $[\ell_i,u_i]$ is the range of $(Tx)_i$ over
$\mathcal P$, computed by linear programming. All sampling parameters are
computed from the base data before the draw, and all have bit length
polynomial in $I$. The algorithm of each theorem is the closure search of
Definition~\ref{def:qp:closure-search} with a fallback on the same draw; it
never rejects or resamples a draw.

\begin{theorem}[Gaussian-like noise on all coefficients]\label{thm:qp:gauss}
Let a base instance with $A\not\succeq0$ be given, and let $(\alpha,T)$, with
$k$ rows, be either the factorization of Lemma~\ref{lem:qp:normalize}
($k=n_-(A)$, $c_{\rm fr}=63/64$) or a supplied convexifying factorization with
frame constant $c_{\rm fr}$. There
is an algorithm that computes from the base data, in time polynomial in $I$,
an integer $b\le\poly(I)$; draws $\gamma_1,\dots,\gamma_n$ independently from
$\mathcal G_{b,\sigma}$; and returns, on every draw, a rational global
minimizer of $F_\gamma$ on $X$ and the optimal value, or reports $X=\emptyset$.
Its expected number of bit operations is at most
$f(n_z)\,C_0^k(1+\mathcal H_{\rm G})(I+1)^{C_1}$, where $C_0,C_1$ are absolute
constants, $f$ is the function of Proposition~\ref{prop:qp:primitives}(b)
($f\equiv1$ if $n_z=0$), and
\[
 \mathcal H_{\rm G}=\prod_{i=1}^k\Bigl[2+c_{\rm fr}^{-1/2}(4k+6)
 +\frac{(1+2k)\alpha(u_i-\ell_i)}{\sigma\sqrt{2\pi}}\Bigr].
\]
For the factorization of Lemma~\ref{lem:qp:normalize},
$\mathcal H_{\rm G}\le\bigl[9+5k+2(1+2k)\nu\diam(\mathcal P)/\sigma\bigr]^k$,
so the expected work is at most
$f_1\bigl(n_z,k,1+\nu\diam(\mathcal P)/\sigma\bigr)(I+1)^{C_1}$ with
$k=n_-(A)$. If $A\succeq0$, one exact solve by
Proposition~\ref{prop:qp:primitives} suffices.
\end{theorem}

The theorem assumes neither growth, nor a unique minimizer, nor genericity,
nor strict complementarity, nor full-dimensional feasibility, nor a supplied
global certificate. The ratio $\nu\diam(\mathcal P)/\sigma$ enters through its
numerical value; a short binary encoding does not make a large ratio
harmless. The exponent $C_1$ does not depend on $k$ or $n_z$, so the bound is
fixed-parameter in $(n_z,k)$ and the numerical ratio. The law
$\mathcal G_{b,\sigma}$ is a finite rational law whose accuracy $b$ depends on
the base instance; it is not an exact Gaussian, and the theorem does not
cover other approximations of a Gaussian.

\begin{theorem}[Uniform noise on all coefficients]\label{thm:qp:uniform}
Let a base instance with $A\not\succeq0$ be given, and let $(\alpha,T)$ be a
convexifying factorization with $k$ rows, $T$ of full row rank, and
$\norm T\le1$, for example that of Lemma~\ref{lem:qp:normalize}. Put
$w_i=u_i-\ell_i+2\sigma\norm{\Psi_{i,:}}_1/\alpha$. There is an algorithm that
computes from the base data a power of two $M$ with $\log M\le\poly(I)$, draws
$\gamma_1,\dots,\gamma_n$ independently from $U_{\sigma,M}$, and
returns an exact minimizer and the optimal value on every draw, with expected
bit work at most $f(n_z)\,C_0^k(1+\mathcal H_{\rm amb})(I+1)^{C_1}$, where
\[
 \mathcal H_{\rm amb}=\prod_{i=1}^k\Bigl[2+\frac{(1+2k)\alpha\sqrt n\,w_i}{2\sigma}\Bigr].
\]
For the factorization of Lemma~\ref{lem:qp:normalize},
$\mathcal H_{\rm amb}\le\bigl[2+(1+2k)\bigl(2n+2\nu\sqrt n\diam(\mathcal P)/\sigma\bigr)\bigr]^k$.
\end{theorem}

This bound is polynomial for each fixed $k$ when the displayed ratios are
polynomially bounded. It is not a fixed-parameter bound: the power of $n$
grows with $k$. Section~\ref{sec:qp:ambient} shows that this power is a real
property of the counts used in the proof.

\begin{theorem}[Noise aligned with the factor]\label{thm:qp:aligned}
Let a continuous base instance ($n_z=0$) and an arbitrary convexifying
factorization $(\alpha,T)$ with $k\ge1$ rows be given, and put
$w_i=u_i-\ell_i+2\sigma/\alpha$. There is an algorithm that computes from the
base data a power of two $M$ with $\log M\le\poly(I)$, draws
$\xi_1,\dots,\xi_k$ independently from $U_{\sigma,M}$, and returns an
exact minimizer of $F(x)+\xi^{\mathsf T} Tx$ and the optimal value on every draw,
with expected bit work at most $C_0^k(1+\mathcal H_{\rm al})(I+1)^{C_1}$, where
\[
 \mathcal H_{\rm al}=\prod_{i=1}^k\Bigl[3+\frac{(1+2k)\alpha w_i}{2\sigma}\Bigr].
\]
\end{theorem}

The aligned bound is fixed-parameter in $k$ and the ratios
$\alpha(u_i-\ell_i)/\sigma$. Its law perturbs only the $k$ factor
coefficients. Mixed-integer instances are excluded because aligned noise does
not perturb the integer coefficients independently, so the label-gap failures
of Lemma~\ref{lem:qp:gap}(c) are not controlled.

\begin{remark}[The original objective]\label{rem:qp:regret}
By Proposition~\ref{prop:model:regret}, an exact minimizer $\hat x$ of
$F_\gamma$ satisfies $F(\hat x)-\min_XF\le\max_{x,y\in X}\gamma^{\mathsf T}(x-y)$,
which is at most $\bar\sigma\sum_iw_i$ when every coefficient is bounded
by $\bar\sigma$; for $\mathcal G_{b,\sigma}$, $\bar\sigma=(b+20)\sigma$.
Under aligned noise the bound is $\bar\sigma\sum_i(u_i-\ell_i)$. Because the
work bounds grow as $\sigma$ decreases, these results do not recover the
unperturbed optimizer, and they claim no improvement of worst-case
approximation bounds for the unperturbed problem.
\end{remark}

\subsection{Aligned, uniform, and Gaussian-like noise}\label{sec:qp:ambient}

The three laws lead to different counts because of the conditional law of
the factor coefficients given the residual. Every local event of the search,
\[
 E_v:\qquad V(v)\le V(v\pm h_{ij}e_i)+2E_j\quad(i\in Q_v),
\]
confines each $\xi_i$ with $i\in Q_v$ to an interval of length at most
$(1+2k)\alpha h_{ij}$ whose position depends on $\zeta$
(Lemma~\ref{lem:count:interval} applied to $W_\zeta$). Under aligned noise
$\zeta=0$ and the factor coefficients are independent, so
Corollary~\ref{cor:count:levels} applies directly. Under uniform ambient
noise, $\xi=\Psi\gamma$ is dependent on $\zeta=\Pi_T\gamma$, and the
residual-dependent intervals must be handled by a volume estimate.

\begin{lemma}[Dependent fiber volume]\label{lem:qp:volume}
Let $\gamma$ be uniform on $[-\sigma,\sigma]^n$, let $T\in\R^{k\times n}$ have
full row rank and $\norm T\le1$, and let $Q\subseteq\{1,\dots,k\}$ have $t$
elements. For each $i\in Q$ let $I_i(\zeta)$ be an interval of length at most
$\ell_i$ depending measurably on $\zeta=\Pi_T\gamma$. Then
\[
 \Prob\{\xi_i\in I_i(\zeta)\ (i\in Q)\}
 \le\frac{\prod_{i\in Q}\ell_i}{(2\sigma)^t}\sum_{|K|=t}|\det T_{Q,K}|
 \le\prod_{i\in Q}\frac{\sqrt n\,\ell_i}{2\sigma}.
\]
\end{lemma}

The factor $\sqrt n$ is attained even for perfectly conditioned factors.

\begin{proposition}[Sharpness of the fiber estimate]\label{prop:qp:fiber-sharp}
Let $n=m^2$, $T=m^{-1}(1,\dots,1)$, and $\gamma$ uniform on
$[-\sigma,\sigma]^n$. Then $TT^{\mathsf T}=1$, $\operatorname{Var}\xi=\sigma^2/3$,
and for the interval $I(\zeta)$ of length $\ell\le2\sigma m$ centered at the
midpoint of the conditional support of $\xi$,
$\Prob\{\xi\in I(\zeta)\}=1-(1-\ell/(2\sigma m))^n$, whose ratio to $\ell$
tends to $\sqrt n/(2\sigma)$ as $\ell\to0$. For $k$ disjoint blocks of this
form, in dimension $N=km^2$, the ratio of the joint probability to
$\prod_i\ell_i/(2\sigma)$ tends to $m^k=(N/k)^{k/2}$.
\end{proposition}

A dimension-free variance of $\xi$ does not imply dimension-free
concentration on intervals that may depend on the residual. The loss is not
confined to arbitrary intervals: it occurs for the counts of actual quadratic
instances.

\begin{proposition}[The uniform-noise counts grow with the dimension]\label{prop:qp:ambient-barrier}
For every $k\ge1$ and every even $m\ge256$ there is a rational continuous
instance with $N=km^2$ variables, $O(N)$ inequalities, $n_-(A)=k$, and a
convexifying factorization with $\alpha=1$, $TT^{\mathsf T}=I_k$, and $\alpha<4\nu$,
whose ranges satisfy $u_i-\ell_i=8$, such that under independent uniform
noise on all coefficients, either continuous on $[-1,1]$ or $U_{1,M}$
with $M\ge32$, the following hold for the meshes on the box of
Theorem~\ref{thm:qp:uniform}. At some level the expected number of nodes
satisfying the local event is at least $128^{-k}(N/k)^{k/2}$, and at some
level the expected number of $2E_j$-nearly optimal nodes is at least
$64^{-k}(N/k)^{k/4}$.
\end{proposition}

The proposition is the case $\alpha=1$ of Theorem~\ref{thm:lim:ambient}, which
is proved in Appendix~\ref{app:lim:ambient} for a slightly more general
family. Consequently no bound of the form $f(k,\mathsf R)N^{C}$ with an
absolute exponent $C$ holds for either count, where $\mathsf R$ records the projected
widths, the frame constant, and the curvature-to-noise ratios. The proposition
limits the counting argument, not the closure algorithm, which may close the
few quadratic pieces of these instances early. The ambient diameter of these
instances grows with $m$, so the proposition does not exclude a bound
parameterized by the full diameter.

Gaussian noise removes both losses. For $\gamma\sim N(0,\sigma^2I_n)$, the
projections $\xi=\Psi\gamma$ and $\zeta=\Pi_T\gamma$ are jointly Gaussian and
uncorrelated, hence independent, and
$\operatorname{Cov}\xi=\sigma^2(TT^{\mathsf T})^{-1}$. Conditioning on $\zeta$ does not
change the law of $\xi$, whose density obeys a product bound that depends only
on the frame constant. In addition, Lemma~\ref{lem:qp:aux}(c) locates the
allowed values of $\xi_i$: combined with the two neighbor comparisons it
gives $|\xi_i+\alpha(v_i-(Tx_v)_i)|\le(1+2k)\alpha h_{ij}/2$, so
$\xi_i\in\alpha[\ell_i-v_i,u_i-v_i]+[-(1+2k)\alpha h_{ij}/2,(1+2k)\alpha h_{ij}/2]$,
an interval determined by $v$ and the projected range rather than by the
enlarged search box.

\begin{lemma}[Gaussian-weighted local counts]\label{lem:qp:gauss-count}
Let $\gamma\sim N(0,\sigma^2I_n)$, let $T$ have frame constant $c_{\rm fr}$,
let $\Lambda=\diag(\lambda_i)$, let $[\ell_i,u_i]$ contain the range of
$(Tx)_i$ over $X$, and use balanced nested meshes on any fixed box with the
constant $C_k$ of Definition~\ref{def:count:mesh} for the curvature vector
$(\lambda_1,\dots,\lambda_k)$. At every level $j$,
\[
 \sum_{v\in G_j}\Prob(E_v)\le\prod_{i=1}^k\Bigl[2+c_{\rm fr}^{-1/2}(2C_k+4)
 +\frac{C_k\lambda_i(u_i-\ell_i)}{\sigma\sqrt{2\pi}}\Bigr].
\]
The bound does not depend on the location or the size of the box.
\end{lemma}

The finite law $\mathcal G_{b,\sigma}$ has none of these independence
properties. Its local-event probabilities are obtained from the Gaussian
proxy by Lemma~\ref{lem:count:transfer}, using a uniform bound on the
number of intervals in every coordinate section of a local event
(Lemma~\ref{lem:qp:sections}); its tube probabilities are transferred in the
same way.

\subsection{Separable mixed recourse with rational breakpoints}\label{sec:qp:separable}

When the convex part is separable and the domain is a product, the recourse
decomposes into scalar problems. Integer winners are then certified globally
by neighbor inequalities, and the number of integer variables no longer
enters the complexity.

\begin{definition}[Separable instances]\label{def:qp:sep}
The domain is $X=\prod_{j=1}^nX_j$, where $X_j=[\ell_j,u_j]$ with
$\ell_j<u_j$ (continuous) or $X_j=\Z\cap[\ell_j,u_j]$ with at least two
points (integer); integer endpoints are rounded inward, empty domains are
rejected, and fixed coordinates are removed beforehand, their contributions
being absorbed into the remaining data. Each $\varphi_j$ is a continuous
convex function on $[\ell_j,u_j]$ given by rational breakpoints
$\ell_j=\kappa_{j0}<\dots<\kappa_{js_j}=u_j$ and rational quadratics
$\psi_{jm}(t)=\frac12c_{jm}t^2+\mu_{jm}t$ (plus constants) with
$c_{jm}\ge0$ on $[\kappa_{j,m-1},\kappa_{jm}]$. The objective is
\[
 F(x)=\sum_{j=1}^n\varphi_j(x_j)-\frac\alpha2\norm{Tx}^2,
 \qquad\alpha>0,\quad T\in\Q^{k\times n}.
\]
The bit length $I$ counts all pieces, breakpoints, and binary endpoints.
Continuity, the curvature signs, and the order of the one-sided derivatives at
the breakpoints are checked exactly.
\end{definition}

Rational breakpoints are essential. The function $\max\{0,x^2-2\}$ on $[0,2]$
has rational quadratic pieces, but $\max\{0,x^2-2\}-x$ has its unique minimizer
at the irrational breakpoint $\sqrt2$. The parameters $k$ and $\alpha$ are
those of the supplied concave term, not the negative inertia or curvature of
the whole objective, and the product domain excludes coupling constraints.

\begin{lemma}[Scalar recourse]\label{lem:qp:scalar}
For $\lambda\in\R$ consider $\min_{t\in X_j}[\varphi_j(t)-\lambda t]$.
\begin{enumerate}[label=(\alph*)]
\item If $X_j$ is integer, put $\delta_j(t)=\varphi_j(t+1)-\varphi_j(t)$, which
 is nondecreasing. A point $t\in X_j$ is optimal if and only if
 $\delta_j(t-1)\le\lambda\le\delta_j(t)$, where a condition involving a point
 outside $X_j$ is omitted. An optimal $t$ is found by binary search with
 $O(\log(u_j-\ell_j+2))$ evaluations of $\varphi_j$.
\item If $X_j$ is continuous, the following at most $2s_j+1$ states cover all
 $\lambda\in\R$, and each certifies global optimality on its closed interval
 of $\lambda$: for a piece with $c_{jm}>0$, the free state
 $t=(\lambda-\mu_{jm})/c_{jm}$ for $\lambda$ between the derivatives of
 $\psi_{jm}$ at $\kappa_{j,m-1}$ and $\kappa_{jm}$; and for every breakpoint
 $\kappa_{jm}$, the constant state $t=\kappa_{jm}$ for $\lambda$ between the
 left and right derivatives of $\varphi_j$ at $\kappa_{jm}$, with $-\infty$
 and $+\infty$ replacing the missing derivatives at $\ell_j$ and $u_j$.
\end{enumerate}
\end{lemma}

With $\lambda_j(a,\zeta)=\alpha t_j^{\mathsf T} a-\zeta_j$, where $t_j$ is column $j$
of $T$, the inner problem defining $W_\zeta(a)$ equals
$\frac\alpha2\norm a^2+\sum_j\min_{x_j\in X_j}[\varphi_j(x_j)-\lambda_j(a,\zeta)x_j]$.
A choice of one state per coordinate gives a piece in the sense of
Lemma~\ref{lem:qp:pieces}, with a region described by at most $2n$
inequalities with normals $\pm\alpha t_j^{\mathsf T}$ and offsets affine in $\zeta$,
on which the piece's response attains $W_\zeta$ itself. There are at most
$R_{\rm sep}=\prod_{j\text{ integer}}N_j\prod_{j\text{ continuous}}(2s_j+1)$
pieces, where $N_j=|X_j|$. A piece has Hessian
$\alpha I-\alpha^2\sum_{j\text{ free}}t_jt_j^{\mathsf T}/c_j$, bounded in norm by
$\Theta_{\rm sep}=\alpha+\alpha^2\sum_{j\text{ continuous}}r_j^+\norm{t_j}^2$,
where $r_j^+$ is the largest reciprocal of a positive listed curvature of
$\varphi_j$ (zero if none). Since every region certifies the global value, no
gap test is needed. The exact fallback enumerates integer assignments and
continuous pieces and minimizes the resulting quadratics on the boxes of
pieces by Lemma~\ref{lem:qp:face}(d); its multiplier is
$B_{\rm sep}=\max\{2,\prod_{j\text{ integer}}N_j\prod_{j\text{ continuous}}3s_j\}$.
All these counts have logarithms polynomial in $I$.

\begin{theorem}[Separable recourse: aligned and uniform noise]\label{thm:qp:sep}
Let a separable instance with $k\ge1$ be given, and let $[\ell_i,u_i]$ be the
range of $(Tx)_i$ over $X$.
\begin{enumerate}[label=(\alph*)]
\item Under $\xi_1,\dots,\xi_k$ drawn independently from $U_{\sigma,M}$,
 with $M$ computed from the base data, there is an algorithm that returns an
 exact rational minimizer of $F(x)+\xi^{\mathsf T} Tx$ and the optimal value on every
 draw, with expected bit work at most $C_0^k(1+\mathcal H_{\rm al})(I+1)^{C_1}$,
 where $\mathcal H_{\rm al}$ is as in Theorem~\ref{thm:qp:aligned}.
\item If $T$ has full row rank and $\norm T\le1$, the same holds under
 $\gamma_1,\dots,\gamma_n$ drawn independently from $U_{\sigma,M}$,
 with the bound of Theorem~\ref{thm:qp:uniform} without the factor $f(n_z)$.
\end{enumerate}
The number of integer variables is not a parameter.
\end{theorem}

A general spectral normalization of the separable objective would destroy
the separability of the convex part. A normalization that only rotates and
rescales the rows of the supplied factor preserves it exactly, at the price
of an anisotropic auxiliary metric.

\begin{lemma}[Anisotropic row normalization]\label{lem:qp:aniso}
Let $T\in\Q^{k\times n}$ have full row rank and put $\beta=\alpha\norm T^2$.
In polynomial time one computes an exactly orthogonal rational $Q$ and dyadic
numbers $d_1,\dots,d_k>0$ such that $U=\diag(d_i)^{-1}QT$ and
$\Lambda=\diag(\alpha d_i^2)$ satisfy
\[
 U^{\mathsf T}\Lambda U=\alpha T^{\mathsf T} T,\qquad
 \tfrac1{16}I_k\preceq UU^{\mathsf T}\preceq I_k,\qquad
 \max_i\lambda_i<8\beta .
\]
\end{lemma}

The identity $U^{\mathsf T}\Lambda U=\alpha T^{\mathsf T} T$ is exact. The approximation in
the Jacobi rotations is used only to bound the frame; it is neither discarded
nor moved into the convex part. Small singular values of $T$ affect only the
bit lengths of the transformed data.

\begin{theorem}[Separable recourse: Gaussian-like noise]\label{thm:qp:sep-gauss}
Let a separable instance with $k\ge1$ and $T$ of full row rank after
preprocessing be given; no norm or conditioning promise on $T$ is needed. Let
$U$ and $\Lambda$ be as in Lemma~\ref{lem:qp:aniso} and let $[\ell_i,u_i]$ be
the range of $(Ux)_i$ over $X$. There is an algorithm that computes from the
base data an integer $b\le\poly(I)$, draws $\gamma_1,\dots,\gamma_n$
independently from $\mathcal G_{b,\sigma}$, and returns an exact rational
minimizer of $F_\gamma$ and the optimal value on every draw, with expected bit
work at most $C_0^k(1+\mathcal H_\Lambda)(I+1)^{C_1}$, where, with
$C_k=1+8k$,
\[
 \mathcal H_\Lambda=\prod_{i=1}^k\Bigl[2+4(2C_k+4)+\frac{C_k\lambda_i(u_i-\ell_i)}{\sigma\sqrt{2\pi}}\Bigr]
 \le\Bigl[8C_k+18+\frac{8C_k\beta\diam(X)}{\sigma\sqrt{2\pi}}\Bigr]^k .
\]
Thus the expected work is $f(k,1+\beta\diam(X)/\sigma)(I+1)^{C_1}$ with an
absolute exponent and no parameter for the number of integer variables.
\end{theorem}

The search uses meshes balanced by the curvatures: dyadic scales $\rho_i$
with $1\le\lambda_i\rho_i^2<4$ give $c_{\rm bal}<4$ in
Definition~\ref{def:count:mesh}, so that no ratio of largest to smallest
curvature enters the count. The parameter $\beta$ is the curvature of the
supplied concave term; it can be much larger than the negative curvature of
the whole objective. The theorem does not identify a minimum-rank
representation or remove dependent rows of $T$.

\subsection{Summary of the regimes}\label{sec:qp:summary}

\begin{table}[ht]
\centering
\small
\begin{tabularx}{\textwidth}{@{}lXXX@{}}
\toprule
Result & Law (fixed before sampling) & Recourse and domain & Expected bit work \\
\midrule
Thm.~\ref{thm:qp:two} & $U_{\sigma,M}$ or $\mathcal G_{b,\sigma}$ on all $n$ coefficients & continuous polytope, $k\le2$ & $\poly(I)(1+\nu S/\sigma)$; linear in the ratio \\
Thm.~\ref{thm:qp:aligned} & $U_{\sigma,M}$ on the $k$ factor coefficients & continuous polytope & $C^k(1+\mathcal H_{\rm al})\poly(I)$; parameters $k$ and $\alpha(u_i-\ell_i)/\sigma$ \\
Thm.~\ref{thm:qp:uniform} & $U_{\sigma,M}$ on all $n$ coefficients & mixed polytope & $f(n_z)C^k(1+\mathcal H_{\rm amb})\poly(I)$; a power of $n$ growing with $k$ \\
Thm.~\ref{thm:qp:gauss} & $\mathcal G_{b,\sigma}$ on all $n$ coefficients & mixed polytope & $f_1(n_z,k,1+\nu\diam(\mathcal P)/\sigma)\poly(I)$ for intrinsic normalization \\
Thm.~\ref{thm:qp:sep} & aligned or uniform & separable, product domain & as Thms.~\ref{thm:qp:aligned}, \ref{thm:qp:uniform}, no $n_z$ \\
Thm.~\ref{thm:qp:sep-gauss} & $\mathcal G_{b,\sigma}$ on all $n$ coefficients & separable, product domain & $f(k,1+\beta\diam(X)/\sigma)\poly(I)$, no $n_z$ \\
\bottomrule
\end{tabularx}
\caption{Exact algorithms for sampled problems with few negative directions.
All are correct on every draw and output rational minimizers and values of
polynomial bit length. In the uniform-noise row, the power of $n$ is a property
of the counting surrogate (Proposition~\ref{prop:qp:ambient-barrier}). For the
intrinsic factorization, $\alpha<4\nu$.}
\label{tab:qp:regimes}
\end{table}

The polynomial exponents are absolute in all rows. The Gaussian-like theorem
is the strongest general statement. The aligned theorem is the simplest and
needs no normalization, but its noise is low-dimensional. The two-direction
theorem has a different and, for $k\le2$, possibly sharper numerical
dependence. The separable theorems trade a restricted convex part for an
unrestricted number of integer variables.
